% 本文件由 tools/build_book.py 从 dist/reports/assertions.json 生成，请勿手改。
\begin{longtable}{@{}p{0.06\textwidth}p{0.32\textwidth}p{0.22\textwidth}p{0.16\textwidth}@{}}
\caption{章节 2.1 数值断言（生成时间 2026-10-18T14:28:42-0400）}\\
\toprule
类别 & 断言 & 理论值 & 实测值 \\
\midrule
\endfirsthead
\toprule
类别 & 断言 & 理论值 & 实测值 \\
\midrule
\endhead
\bottomrule
\endlastfoot

A & 恒等式残差/(u$\cdot$$\Lambda$\_\allowbreak{}n) / 朴素 $\Sigma$ w/(x$-$xi) 的最小值 & 正确恒等式含 $\prod$(\allowbreak{}x$-$x\_j) 因子(\allowbreak{}残差 $\approx$ u$\cdot$$\Lambda$\_n)；朴素和 =\allowbreak{} 1/\allowbreak{}$\prod$(\allowbreak{}x$-$x\_j) 而不是 1 & 归一化残差 $\le$ 5.93(\allowbreak{}equi n=\allowbreak{}6: 比值 1.32；equi n=\allowbreak{}12: 比值 1.10；equi n=\allowbreak{}20: 比值 0.83；cheb n=\allowbreak{}6: 比值 3.84；cheb n=\allowbreak{}12: 比值 3.94；cheb n=\allowbreak{}20: 比值 5.93)；朴素和最小 22.8(\allowbreak{}n=\allowbreak{}6/\allowbreak{}12/\allowbreak{}20 分别为 22.8/\allowbreak{}4373.4/\allowbreak{}2065.6) \checkmark \\
A & max|$\Sigma$$\ell$$-$1| / 同批节点的 max $\Lambda$\_\allowbreak{}n & $\Sigma$$\ell$ $\equiv$ 1(\allowbreak{}舍入级)；$\Lambda$\_n =\allowbreak{} $\Sigma$|$\ell$| 可以任意大------绝对值的位置决定一切 & 单位分解偏差 $\le$ 4.85e-13；$\Lambda$ 最大 10772.4(\allowbreak{}等距 n=\allowbreak{}20 为 2.9) \checkmark \\
A & 乘积形式权重 / 闭式权重 的常数性与散布 & w\_j =\allowbreak{} (\allowbreak{}$-$1)\textasciicircum{}j $\delta$\_j 与 $\prod$ 形式只差整体常数(\allowbreak{}缩放不改变重心插值) & n=\allowbreak{}6: 比值 5.333333，散布 8.3e-16；n=\allowbreak{}12: 比值 170.666667，散布 3.2e-15；n=\allowbreak{}20: 比值 26214.400000，散布 7.9e-15 \checkmark \\
A & exact\_\allowbreak{}solve vs Bareiss 的相对偏差(6 组) & 两种独立的精确有理算法必须给出同一组系数 & 最大 0.0e+00；equi n=\allowbreak{}6: 0.0e+00(\allowbreak{}$\kappa$=\allowbreak{}4.6e+02)；equi n=\allowbreak{}10: 0.0e+00(\allowbreak{}$\kappa$=\allowbreak{}4.0e+04)；equi n=\allowbreak{}12: 0.0e+00(\allowbreak{}$\kappa$=\allowbreak{}3.8e+05)；cheb n=\allowbreak{}6: 0.0e+00(\allowbreak{}$\kappa$=\allowbreak{}2.0e+02)；cheb n=\allowbreak{}10: 0.0e+00(\allowbreak{}$\kappa$=\allowbreak{}6.3e+03)；cheb n=\allowbreak{}12: 0.0e+00(\allowbreak{}$\kappa$=\allowbreak{}3.6e+04) \checkmark \\
A & log2 $\kappa$(V)/n 与 Chebyshev 理论底数 / n=30 的 $\kappa$ 之比 & Chebyshev 节点 $\kappa$ \textasciitilde{} (\allowbreak{}1+$\surd$2)\textasciicircum{}n；等距节点更快；同 n 下 Chebyshev 小几个数量级 & 等距 log2$\kappa$/\allowbreak{}n 1.5451--1.6050(\allowbreak{}漂移，底数 > 1+$\surd$2)；Chebyshev log2$\kappa$/\allowbreak{}n 1.2614--1.2667(\allowbreak{}理论 1.2716，偏差 8.0e-03)；$\kappa$\_eq/\allowbreak{}$\kappa$\_ch(\allowbreak{}n=\allowbreak{}30)=\allowbreak{}965.5 \checkmark \\
A & 三条误差通道(n=36，两种分布) & 函数值误差两条通道同量级；系数误差被 $\kappa$u 界住；残差 $\approx$ u 与 $\kappa$ 无关 & equi n=\allowbreak{}36: 函数值 1.05e-07/\allowbreak{}3.52e-08(\allowbreak{}比 2.98$\times$)，系数 1.19e+00 $\le$ $\kappa$u=\allowbreak{}2.55e+01，残差 1.5e-17；cheb n=\allowbreak{}36: 函数值 1.33e-15/\allowbreak{}8.88e-16(\allowbreak{}比 1.50$\times$)，系数 1.31e-04 $\le$ $\kappa$u=\allowbreak{}5.87e-03，残差 2.2e-17 \checkmark \\
B & |$\ell$\_\allowbreak{}i(x\_\allowbreak{}j)$-$$\delta$\_\allowbreak{}ij| / |$\Sigma$w|/$\Sigma$|w| / |p(x\_\allowbreak{}i)$-$y\_\allowbreak{}i| & Lagrange 基在节点上是单位阵；乘积形式权重之和恒为 0(\allowbreak{}相对残差 $\approx$ u)；插值在节点上取给定值 & 基 0.00e+00；|$\Sigma$w|/\allowbreak{}$\Sigma$|w| 4.16e-17(\allowbreak{}等距 n=\allowbreak{}12 时 1.09e-16，$\Sigma$|w|=\allowbreak{}1.024e+03)；节点值 0.00e+00 \checkmark \\
B & 缩放不变性 / 置换不变性(以 u$\cdot$$\Lambda$\_\allowbreak{}n 归一) & 重心公式分子分母同阶齐次；插值多项式与节点编号无关 & 缩放 1.34u$\Lambda$、\allowbreak{}置换 1.97u$\Lambda$(\allowbreak{}$\Lambda$=\allowbreak{}89.3) \checkmark \\
B & 帧内 $\kappa$ 单调不减 / 后向残差 / 三条通道齐备 & $\kappa$ 随 n 增长；残差与 $\kappa$ 无关(\allowbreak{}后向稳定)；帧必须同时给出三条误差通道 & t=\allowbreak{}6: $\kappa$=\allowbreak{}4.55e+02,\allowbreak{} 残差=\allowbreak{}5.8e-17；t=\allowbreak{}12: $\kappa$=\allowbreak{}3.81e+05,\allowbreak{} 残差=\allowbreak{}5.8e-17；t=\allowbreak{}20: $\kappa$=\allowbreak{}3.19e+09,\allowbreak{} 残差=\allowbreak{}5.8e-17；t=\allowbreak{}30: $\kappa$=\allowbreak{}2.59e+14,\allowbreak{} 残差=\allowbreak{}5.8e-17；残差最大 5.8e-17 \checkmark \\
B & 四支 viz 的帧契约 & requiredAux 齐备、\allowbreak{}数组与 meta 全有限、\allowbreak{}对数轴登记地板、\allowbreak{}pitfalls $\ge$ 60 字 & runge\_nodes: 17 数组 /\allowbreak{} 27 meta /\allowbreak{} 地板键 ['error\_floor',\allowbreak{} 'n\_error\_at\_floor'] /\allowbreak{} pitfalls 359 字；lebesgue\_growth: 17 数组 /\allowbreak{} 25 meta /\allowbreak{} 地板键 ['log10\_floor',\allowbreak{} 'n\_lambda\_at\_floor'] /\allowbreak{} pitfalls 400 字；error\_curve: 15 数组 /\allowbreak{} 27 meta /\allowbreak{} 地板键 ['error\_floor',\allowbreak{} 'n\_error\_at\_floor',\allowbreak{} 'n\_monic\_at\_floor'] /\allowbreak{} pitfalls 410 字；barycentric\_vs\_vandermonde: 17 数组 /\allowbreak{} 24 meta /\allowbreak{} 地板键 ['log10\_floor',\allowbreak{} 'n\_error\_at\_floor'] /\allowbreak{} pitfalls 496 字 \checkmark \\
A & 三种插值形式的最大偏差 / (u$\cdot$$\Lambda$\_\allowbreak{}n) & 同一多项式，差异应落在 c$\cdot$u$\cdot$$\Lambda$\_n(\allowbreak{}c 为小常数) & 最大 10.722(\allowbreak{}equi n=\allowbreak{}8: 1.46/\allowbreak{}1.10；equi n=\allowbreak{}16: 1.38/\allowbreak{}1.00；equi n=\allowbreak{}24: 1.77/\allowbreak{}1.60；cheb n=\allowbreak{}8: 7.03/\allowbreak{}3.52；cheb n=\allowbreak{}16: 5.87/\allowbreak{}4.40；cheb n=\allowbreak{}24: 10.72/\allowbreak{}2.68) \checkmark \\
B & 误差/包络的最大值(严格上界不允许超过 1) & $\le$ 1 & 0.45398 \checkmark \\
A & 包络紧度(误差/包络的中位数) & 严格上界 $\ge$ 实测误差；对 e\textasciicircum{}x 应只宽松一个常数倍 & 中位数 0.360--0.381 \checkmark \\
A & max|$\prod$|(Chebyshev)/ 2\textasciicircum{}(1$-$n) & 首一 Chebyshev 多项式的无穷范数恰为 2\textasciicircum{}(\allowbreak{}1$-$n)：网格最大值只能 $\le$ 理论值 & 0.9285--0.9969(\allowbreak{}n=\allowbreak{}4 网格最粗，n=\allowbreak{}20 最接近 1) \checkmark \\
A & max|$\prod$|(等距)/ max|$\prod$|(Chebyshev) & 等距节点的首一乘积在端点附近爆炸，比值随 n 增长(\allowbreak{}n=\allowbreak{}12 实测 8.3) & 0.978--122.889 \checkmark \\
A & 小 n 的 $\Lambda$\_\allowbreak{}n 相对误差(8 个精确值) & $\Lambda$\_1=\allowbreak{}1、\allowbreak{}$\Lambda$\_2=\allowbreak{}1.25、\allowbreak{}$\Lambda$\_3(\allowbreak{}ch)=\allowbreak{}5/\allowbreak{}3、\allowbreak{}$\Lambda$\_4(\allowbreak{}ch)=\allowbreak{}1.7987618 等 & 最大相对误差 2.507e-11 \checkmark \\
A & Chebyshev 渐近偏差 / 等距渐近比值 & Chebyshev: $\Lambda$\_n $\approx$ (\allowbreak{}2/\allowbreak{}$\pi$)ln n + 0.9625；等距: $\Lambda$\_n $\approx$ 2\textasciicircum{}(\allowbreak{}n+1)/\allowbreak{}(\allowbreak{}e n ln n) & Chebyshev 偏差 $\le$ 3.05e-03；等距比值 0.853--0.914(\allowbreak{}n$\ge$15 时 0.853--0.863) \checkmark \\
A & $\Lambda$\_\allowbreak{}\{n+1\}/$\Lambda$\_\allowbreak{}n(等距)与两侧斜率偏差 & 等距 $\Lambda$ \textasciitilde{} 2\textasciicircum{}n(\allowbreak{}斜率 =\allowbreak{} n$\cdot$ln2/\allowbreak{}ln10)；Chebyshev $\Lambda$ \textasciitilde{} (\allowbreak{}2/\allowbreak{}$\pi$)ln n & $\Lambda$\_30/\allowbreak{}$\Lambda$\_29=\allowbreak{}1.9146；等距斜率偏差 $\le$ 0.0183(\allowbreak{}n$\ge$20；n=\allowbreak{}10/\allowbreak{}15 还有 0.220)；Chebyshev 斜率偏差 $\le$ 0.0028(\allowbreak{}n=\allowbreak{}10..25) \checkmark \\
A & $\Lambda$\_\allowbreak{}30 等距 / Chebyshev & n=\allowbreak{}30 时等距 6.60e6、\allowbreak{}Chebyshev 3.127，量级差 6 个数量级 & 6.6011e+06 /\allowbreak{} 3.126968 \checkmark \\
A & 等距误差最小值 / 奇偶两条子列的单调性 & 等距插值在 n=\allowbreak{}7 取到全局最小；误差**不是**随 n 单调------按奇偶分成两条子列各自发散 & n=\allowbreak{}7，min=\allowbreak{}0.2474；奇列 n=\allowbreak{}7$\to$29：0.247$\to$333.0，偶列 n=\allowbreak{}8$\to$30：1.045$\to$2388.0 \checkmark \\
A & Chebyshev 误差(n=10/20/30)与奇点距离 & Chebyshev 节点下解析函数一致收敛；奇点距离 1/\allowbreak{}$\surd$a 决定收敛速率 & 1.3219e-01 /\allowbreak{} 1.7733e-02 /\allowbreak{} 2.4256e-03；1/\allowbreak{}$\surd$a=\allowbreak{}0.200，$\rho$=\allowbreak{}1.2198；n=\allowbreak{}30 奇列 6.17e-03 /\allowbreak{} 偶列 2.43e-03 \checkmark \\
A & Remez 交错点数 / 振幅比 / 下界与 max|e| 的相对偏差 & 等幅振荡点数 $\ge$ m+2；min|e(\allowbreak{}x\_i)| $\le$ E =\allowbreak{} max|e|(\allowbreak{}交错定理的两个方向) & exp m=\allowbreak{}4: n\_ext=\allowbreak{}6、\allowbreak{}振幅比 1.000000、\allowbreak{}下界/\allowbreak{}E=\allowbreak{}1.000000；exp m=\allowbreak{}8: n\_ext=\allowbreak{}10、\allowbreak{}振幅比 1.000000、\allowbreak{}下界/\allowbreak{}E=\allowbreak{}1.000000；abs m=\allowbreak{}5: n\_ext=\allowbreak{}7、\allowbreak{}振幅比 1.000000、\allowbreak{}下界/\allowbreak{}E=\allowbreak{}1.000000；sin m=\allowbreak{}6: n\_ext=\allowbreak{}8、\allowbreak{}振幅比 1.000000、\allowbreak{}下界/\allowbreak{}E=\allowbreak{}1.000000；runge m=\allowbreak{}6: n\_ext=\allowbreak{}9、\allowbreak{}振幅比 1.000000、\allowbreak{}下界/\allowbreak{}E=\allowbreak{}1.000000 \checkmark \\
A & E\_\allowbreak{}remez vs Chebyshev 插值误差(m=4, e\textasciicircum{}x) & 最佳一致逼近不差于任何同次数多项式，特别是 Chebyshev 插值 & E\_remez=\allowbreak{}5.466676e-04 $\le$ E\_cheb=\allowbreak{}1.065941e-03 \checkmark \\
A & 重心公式误差 / (u$\cdot$$\Lambda$\_\allowbreak{}n)(n=16..30) & 插值误差已在舍入底以下时，实测误差 $\approx$ c$\cdot$u$\cdot$$\Lambda$\_n(\allowbreak{}c 为个位数常数) & 1.130--4.404 \checkmark \\
A & n=10 时的同一比值(插值误差仍占主导) & 插值误差 $\gg$ u$\cdot$$\Lambda$\_n 时，这个比值不再反映舍入行为 & 等距 7.082e+04、\allowbreak{}Chebyshev 1.857e+05 \checkmark \\
A & Vandermonde 系数误差 $\le$ $\kappa$(V)$\cdot$u & 前向误差界 $\|$$\delta$c$\|$/\allowbreak{}$\|$c$\|$ 不超过 $\kappa$(\allowbreak{}V)$\cdot$u 的常数倍(\allowbreak{}最坏情况尺度) & equi n=\allowbreak{}20: 4.18e-09 $\le$ 3.54e-07；equi n=\allowbreak{}30: 2.69e-04 $\le$ 2.88e-02；cheb n=\allowbreak{}20: 8.41e-11 $\le$ 4.50e-09；cheb n=\allowbreak{}30: 1.14e-06 $\le$ 2.98e-05 \checkmark \\
A & Vandermonde 后向残差(n=6..30 的最大值) & 部分主元 LU 后向稳定：残差 $\approx$ u，与 $\kappa$ 无关 & 2.634e-17 \checkmark \\
A & 闭式 vs Cauchy 积分的相对偏差(k=1,2,4,8，r=0.15) & 两条独立路径必须给出同一组 f\textasciicircum{}(\allowbreak{}k)/\allowbreak{}k!(\allowbreak{}奇点在 $\pm$0.2i，故 r<0.2) & k=\allowbreak{}1: 1.49e-16、\allowbreak{}k=\allowbreak{}2: 1.07e-16、\allowbreak{}k=\allowbreak{}4: 2.27e-16、\allowbreak{}k=\allowbreak{}8: 1.49e-16 \checkmark \\
A & 轮廓半径 0.25 > 奇点距离 0.2 时的相对偏差(k=1,4) & 围道内必须解析：圈住极点后积分值被残差主导，函数**不会报错** & k=\allowbreak{}1: 9.99e-01、\allowbreak{}k=\allowbreak{}4: 1.00e+00 \checkmark \\
B & |p(x\_\allowbreak{}i) $-$ y\_\allowbreak{}i| / max|y|(三种实现 $\times$ 两种分布 $\times$ 3 个 n) & p(\allowbreak{}x\_i) =\allowbreak{} y\_i 精确成立(\allowbreak{}浮点下 $\le$ c$\cdot$n$\cdot$u) & 1.634e-16 \checkmark \\
B & |T\_\allowbreak{}n(递推) $-$ cos(n$\cdot$arccos x)| 的最大值(n $\le$ 20) & 递推与三角定义是同一个多项式 & 6.772e-15 \checkmark \\
B & T\_\allowbreak{}4 的单项式系数 & T\_4 =\allowbreak{} 8x\textasciicircum{}4 $-$ 8x\textasciicircum{}2 + 1；首项系数 2\textasciicircum{}(\allowbreak{}n$-$1) & [ 8.  0. -8.  0.  1.](\allowbreak{}T\_3 首项 4) \checkmark \\
B & 实测 sup|T\_\allowbreak{}n/2\textasciicircum{}(n$-$1)| vs 2\textasciicircum{}(1$-$n) & 首一 Chebyshev 多项式的无穷范数恰为 2\textasciicircum{}(\allowbreak{}1$-$n) & n=\allowbreak{}1: 1/\allowbreak{}1；n=\allowbreak{}2: 0.5/\allowbreak{}0.5；n=\allowbreak{}4: 0.125/\allowbreak{}0.125；n=\allowbreak{}8: 0.0078125/\allowbreak{}0.0078125；n=\allowbreak{}12: 0.000488281/\allowbreak{}0.000488281 \checkmark \\
B & 差商表对角线偏差 / 换序后的多项式偏差 & 对角线即差商；插值多项式与节点编号无关(\allowbreak{}差商对节点对称) & 对角线 0.000e+00，换序 1.554e-15 \checkmark \\
B & min\_\allowbreak{}x $\Lambda$\_\allowbreak{}n(x) / max\_\allowbreak{}i |$\Lambda$\_\allowbreak{}n(x\_\allowbreak{}i) $-$ 1| & $\Lambda$\_n(\allowbreak{}x) $\ge$ 1(\allowbreak{}插值算子范数下界)；节点处 $\ell$\_i(\allowbreak{}x\_j)=\allowbreak{}$\delta$\_ij 故 $\Lambda$\_n(\allowbreak{}x\_i)=\allowbreak{}1 & 下界 1.004975；节点偏差 0.000e+00 \checkmark \\
B & 夹取后的取值范围 / 地板计数 / 非有限计数 / 负数是否报错 & 0 与 inf/\allowbreak{}nan 都必须如实计数；负数在对数轴上无定义，必须报错 & 范围 [1.11e-16,\allowbreak{} 1.00e+12]，地板 3.0，非有限 2.0，负数报错=\allowbreak{}True \checkmark \\
B & poly.py 里被禁的 import & P5：numcore 不得依赖 scipy/\allowbreak{}sympy/\allowbreak{}mpmath；P10：numcore 不得 import numcore.reference & 无 \checkmark \\
\end{longtable}
